% Vietnamese translation for OI Public Library
% Source-File: IOI中国国家候选队论文/2006/高逸涵/高逸涵.ppt
\documentclass[11pt]{article}
\usepackage{oipl-vn}

\title{Factory và các truy hồi theo tỉ lệ\\{\large Ghi chú theo slide}}
\author{Gao Yihan}
\date{2006}

\begin{document}
\maketitle

\origtitle{Factory presentation notes}
\sourcepath{IOI China candidate team papers / 2006 / Gao Yihan / Gao Yihan.ppt}

\section{Phạm vi}

Không thấy bản DOC cùng tên trong thư mục 2006. Văn bản trích từ PowerPoint
cho thấy tiêu đề hoặc tên bài \emph{Factory}, một ảnh nền cục bộ, và nhiều
công thức truy hồi với các tham số \(M,N,K\). Vì thiếu thuyết minh đầy đủ,
bản ghi này chỉ chuyển các mốc đọc được thành ghi chú companion.

\section{Các công thức đọc được}

Nhóm công thức đầu có dạng
\[
  A_1=\frac{M}{N}(N-1),\qquad
  A_i=\frac{A_{i-1}}{N}(N-1).
\]
Một nhóm khác trừ đi một lượng trước khi nhân tỉ lệ:
\[
  A_i=\frac{A_{i-1}-K}{N}(N-1).
\]
Slide cũng biến đổi bằng cách đặt \(B_i=A_i+N-1\), đưa truy hồi về dạng nhân
tỉ lệ đều hơn:
\[
  B_i=\frac{B_{i-1}}{N}(N-1).
\]

\section{Cách hiểu thận trọng}

Các mảnh này gợi ý một bài toán trong đó mỗi bước giữ lại tỉ lệ
\((N-1)/N\) của một đại lượng, đôi khi sau khi trừ một chi phí cố định hoặc
chi phí phụ thuộc bước. Phép đổi biến \(B_i=A_i+N-1\) là kỹ thuật quen thuộc
để khử hằng số trong truy hồi tuyến tính bậc nhất.

\section{Ghi chú sử dụng}

Khi chưa có phần phát biểu bài và hình minh họa, không nên suy rộng thành một
lời giải cụ thể. Nội dung đáng tin cậy ở đây là dạng truy hồi, phép đổi biến,
và định hướng phân tích đại lượng còn lại sau nhiều vòng xử lý.

\end{document}
